% 付録再編草案
% 既存の付録を置き換えず,構成と論の流れを検討するための草案.
% 中心に置く問い: 漸化式を「状態の更新」と見ると,何が見えるか.

\section{漸化式を繰り返すと,何が見えるか}

まず,一つの簡単な漸化式を実際に追ってみよう.
\[
  a_0=0,\qquad a_{n+1}=\frac12a_n+1
\]
とする.
最初の項を計算すると,
\[
  0,\quad 1,\quad \frac32,\quad \frac74,\quad \frac{15}{8},\quad\ldots
\]
となる.
項は $2$ に近づいているように見える.
なぜだろうか.

この漸化式の右辺に $2$ を入れると,$\frac12\cdot2+1=2$ となる.
つまり,$2$ はこの更新規則を一度適用しても変わらない値である.
このような値を不動点という.
ただし,不動点を見つけただけで,数列がそこへ近づくと分かったわけではない.
そこで,各項が $2$ からどれだけ離れているかを調べる.
\[
  e_n=2-a_n
\]
と置くと,
\[
  e_{n+1}
  =2-a_{n+1}
  =2-\left(\frac12a_n+1\right)
  =\frac12(2-a_n)
  =\frac12e_n
\]
である.
誤差は一歩進むごとに半分になる.
初めは $e_0=2$ だから $e_n=2\cdot(1/2)^n$ であり,
\[
  a_n=2-2\left(\frac12\right)^n
\]
となる.
ここでは,項そのものではなく,不動点からのずれを見たことで,近づく理由が等比数列に帰着した.

この例で起きたことを,一般の更新規則
\[
  a_{n+1}=f(a_n)
\]
で考えてみる.
もし数列がある値 $\alpha$ に近づき,しかも $f$ がその近くで連続なら,漸化式の両辺で極限をとることができる.
その結果
\[
  \alpha=f(\alpha)
\]
を得る.
したがって,数列が近づく先は不動点でなければならない.
これは「不動点があれば必ずそこへ近づく」という主張ではない.
先ほどの例では,誤差が半分ずつになったことを別に確かめて,実際に近づくと分かった.

不動点の近くで同じ考え方を使うと,より一般的な見通しが得られる.
$a_n=\alpha+e_n$ と置けば,
\[
  e_{n+1}=f(\alpha+e_n)-f(\alpha)
\]
である.
$f$ が微分可能で,$e_n$ が小さいときには
\[
  f(\alpha+e_n)\approx f(\alpha)+f'(\alpha)e_n
\]
だから,
\[
  e_{n+1}\approx f'(\alpha)e_n
\]
となる.
先ほどの例では $f'(\alpha)=1/2$ だった.
一般にも $|f'(\alpha)|<1$ なら,不動点の十分近くでは誤差が縮む.
非線形な漸化式でも,不動点の近くを拡大して見ると,一次近似による線形な更新が現れるのである.

ここで,誤差を調べるために二つの項だけを使ったことに注目しよう.
もし次の項を作るのに,直前の二項が必要なら,二つをまとめて一つの「状態」と考えればよい.
二階線形漸化式
\[
  a_{n+2}=p a_{n+1}+q a_n
\]
に対して
\[
  \boldsymbol v_n=
  \begin{pmatrix}a_{n+1}\\a_n\end{pmatrix}
\]
と置くと,
\[
  \boldsymbol v_{n+1}
  =
  \begin{pmatrix}p&q\\1&0\end{pmatrix}\boldsymbol v_n
\]
となる.
一つ前の項を含む二次元の状態が,毎回同じ行列で更新される.

この行列の固有ベクトルの方向では,状態は一歩ごとにある倍率 $\lambda$ 倍される.
その倍率が存在する条件
\[
  \det\begin{pmatrix}p-\lambda&q\\1&-\lambda\end{pmatrix}=0
\]
を計算すると,
\[
  \lambda^2-p\lambda-q=0
\]
となる.
これは,本文で一般項を求めるために使った特性方程式である.
つまり特性方程式は,二項をまとめた状態の更新を,単純な倍率ごとに分けて読むための方程式だった.
重解や複素根が現れたときの振る舞いも,行列の反復の姿として理解できる.

複素根の意味を,具体的な漸化式で見てみよう.
\[
  T_0(x)=1,\quad T_1(x)=x,\qquad
  T_{n+1}(x)=2xT_n(x)-T_{n-1}(x)
\]
で定まる多項式列を考える.
$x$ を固定すれば,これは先ほどの二階線形漸化式であり,特性方程式は
\[
  r^2-2xr+1=0
\]
である.
特に $x=\cos\theta$ と置くと,この方程式の根は
$\cos\theta\pm i\sin\theta$ となる.
一方,余弦の加法定理から
\[
  \cos((n+1)\theta)
  =2\cos\theta\cos(n\theta)-\cos((n-1)\theta)
\]
が成り立つ.
したがって,この漸化式と初期値を共有する $\cos(n\theta)$ は
\[
  T_n(\cos\theta)=\cos(n\theta)
\]
を満たす.
複素根の偏角が数列の振動を表し,その振動が多項式の公式として書き直された.
同じ更新規則を,数列として見れば特性根の問題に,多項式として見ればチェビシェフ多項式の問題に読むことができる.

\section{一歩の変化を細かくすると,微積分になるか}

先ほどは,不動点からのずれが一歩ごとにどう変わるかを調べた.
そこで比べたのは,更新の前後の差である.
一般の数列でも,隣り合う項の差
\[
  a_{n+1}-a_n
\]
を見れば,一歩でどれだけ変化したかが分かる.
この「差を調べる」考え方は,微分の考え方に似ている.
では,一歩の幅そのものを小さくしたら,何が起こるだろうか.

時刻 $t_n=nh$ における値を $a_n$ とし,一歩の時間を $h$ とする.
一単位時間あたりの変化は
\[
  \frac{a_{n+1}-a_n}{h}
\]
で測れる.
$h$ を小さくしていくと,差分商は微分係数に近づく.
この考えを逆向きに使えば,微分方程式を数列で近似できる.
微分方程式 $y'=F(t,y)$ では,一歩の間の変化をおよそ $hF(t_n,a_n)$ と見積もり,
\[
  a_{n+1}=a_n+hF(t_n,a_n)
\]
と更新する.
これがオイラー法である.
漸化式を通して微分方程式を計算し,一歩を細かくする極限を通して連続的な変化を考える.
差分と微分は,離散と連続をつなぐ二つの見方なのである.

差分を足し合わせると,途中の項が消えて
\[
  a_n-a_0=\sum_{k=0}^{n-1}(a_{k+1}-a_k)
\]
となる.
左辺は最初から最後までの全体の変化であり,右辺は一歩ずつの変化を足したものである.
これは,積分が小さな変化を積み重ねて全体の変化を表すことに対応している.
微分と積分の関係に似た構造が,差分と和の間にもある.

ここまででは,与えられた変化の規則を数列として計算した.
では逆に,ある値へ近づくように更新規則そのものを作ることはできるだろうか.
方程式 $g(x)=0$ の根を求めたいとする.
現在の近似値を $x_n$ とし,その近くで関数のグラフを接線に置き換える.
小さな移動量 $h$ に対して
\[
  g(x_n+h)\approx g(x_n)+g'(x_n)h
\]
だから,接線上で関数の値が $0$ になるように $h$ を選べば,
\[
  h=-\frac{g(x_n)}{g'(x_n)}
\]
となる.
そこで次の近似値を
\[
  x_{n+1}=x_n-\frac{g(x_n)}{g'(x_n)}
\]
と定める.
接線によって次の項を作るこの漸化式が,ニュートン法である.
方程式を解く計算も,根に近づく更新を設計する問題として見直せる.

接線を使うと,なぜ速く近づくのだろうか.
根を $\alpha$ とし,誤差を $e_n=x_n-\alpha$ とする.
根が単純で関数が十分滑らかなら,テイラー展開により
\[
  g(\alpha+e_n)
  =g'(\alpha)e_n+\frac12g''(\alpha)e_n^2+\cdots
\]
と書ける.
この式と導関数の近似をニュートン法の式に代入すると,誤差の一次の項が打ち消され,
\[
  e_{n+1}
  =
  \frac{g''(\alpha)}{2g'(\alpha)}e_n^2
  +O(e_n^3)
\]
が得られる.
通常の更新では誤差はおおよそ定数倍に縮むが,ニュートン法では十分近くで誤差がおおよそ二乗になる.
一次近似を使って次の項を決め,より高い精度の展開によってその誤差を調べる.
テイラー展開は,ニュートン法の収束の速さを説明するために自然に現れるのである.

\section{数列全体を一つの対象にまとめると,何が見えるか}

ここまでは,項から次の項を作る規則や,隣り合う項の差を調べてきた.
しかし,漸化式を解くときに知りたいのは,一つの項だけではない.
数列全体にどんな構造があるのかを,一つの対象として表せないだろうか.

数列の最初の数項を,形式的に
\[
  a_0+a_1x+a_2x^2+a_3x^3+\cdots
\]
と並べてみる.
ここでは $x$ は,項の位置を記録するための印のような役割を果たす.
この無限和を
\[
  A(x)=\sum_{n=0}^{\infty}a_nx^n
\]
と書いたものが母関数である.
まずは収束を気にせず,係数を整理するための記号として扱おう.
母関数の力は,数列を一つの式にまとめることだけでなく,数列に対する操作を式の計算へ翻訳するところにある.

たとえば $A(x)$ に $x$ を掛けると
\[
  xA(x)=a_0x+a_1x^2+a_2x^3+\cdots
\]
となる.
係数はそのままに,各項を一つずつ右へずらした形である.
一方,隣り合う項の積を足し合わせる畳み込み
\[
  c_n=\sum_{k=0}^{n}a_kb_{n-k}
\]
は,対応する母関数どうしを掛けたときの $x^n$ の係数になる.
項をずらす操作は $x$ 倍に,畳み込みは積に移る.
数列で複雑だった操作が,式の世界では見慣れた代数計算になる.

この翻訳をフィボナッチ数列で試そう.
\[
  F_0=0,\quad F_1=1,\quad F_{n+2}=F_{n+1}+F_n
\]
とし,$F(x)=\sum_{n=0}^{\infty}F_nx^n$ と置く.
係数を一つずつ比べると,漸化式は
\[
  F(x)=x+xF(x)+x^2F(x)
\]
という関係にまとまり,
\[
  F(x)=\frac{x}{1-x-x^2}
\]
を得る.
数列を作っていた更新規則が,母関数では分母に姿を現した.
この分母 $1-x-x^2$ は,特性方程式 $r^2-r-1=0$ と偶然似ているのではない.
分母の根は特性根の逆数である.
行列の固有値,特性方程式,母関数の分母は,フィボナッチ数列の同じ線形構造を別の仕方で表している.

母関数は,線形漸化式だけに使えるわけではない.
カタラン数のように,次の項がそれ以前の項どうしの積の和で決まる数列を考える.
\[
  C_0=1,\qquad C_{n+1}=\sum_{k=0}^{n}C_kC_{n-k}
\]
である.
ここで $C(x)=\sum_{n=0}^{\infty}C_nx^n$ と置くと,右辺の和は母関数の積に対応し,
\[
  C(x)=1+xC(x)^2
\]
となる.
項を一つずつ求める漸化式が,母関数についての二次方程式に変わった.
母関数は,線形の場合には特性方程式と結びつき,非線形の場合には積や組合せの構造を式へ移す.

三つの問いをたどると,同じ対象を何度も見直したことに気づく.
更新を繰り返すと不動点や行列が現れ,一歩を細かくすると微積分が現れ,数列全体をまとめると母関数が現れる.
ここで出てきた数学は,別々の話題を並べたものではない.
漸化式を「状態を更新する規則」として理解しようとしたとき,その構造を見えるようにするために現れた言葉である.
一般項を求めることに加えて,更新が何を保ち,何を変え,どこへ進むのかを問うことも,漸化式を理解する方法なのである.
